
%(BEGIN_QUESTION)
% Copyright 2006, Tony R. Kuphaldt, released under the Creative Commons Attribution License (v 1.0)
% This means you may do almost anything with this work of mine, so long as you give me proper credit

Klassifiser følgende stoffer som {\it partikkel}, {\it atom}, {\it grunnstoff}, {\it molekyl}, {\it forbindelse} og/eller {\it blanding}. Merk at noen av stoffene kan beskrives korrekt med mer enn én klassifisering:

\begin{itemize}
\item{} Ett gram helt rent vann
\item{} Et proton
\item{} Oksygen
\item{} Karbondioksid (CO$_{2}$)
\item{} En kvark
\item{} Bensin
\item{} Luft
\item{} Plutonium
\end{itemize} 

\vskip 20pt \vbox{\hrule \hbox{\strut \vrule{} {\bf Forslag til sokratisk diskusjon} \vrule} \hrule}

\begin{itemize}
\item{} Hvorfor bryr vi oss om disse begrepene? Hvilken praktisk betydning har det for eksempel å vite om karbondioksid er et grunnstoff eller en forbindelse?
\end{itemize}

\underbar{file i00551}
%(END_QUESTION)





%(BEGIN_ANSWER)

Merk: avsnittet ``Terms and Definitions'' i kapittelet ``Chemistry'' i læreboken {\it Lessons In Industrial Instrumentation} er en god ressurs for å svare på disse spørsmålene.

\vskip 10pt

En {\bf partikkel} er en del av et atom, og kan skilles fra andre deler bare ved energinivåer langt over vanlige kjemiske reaksjoner.

\vskip 10pt

Et {\bf atom} er den minste enheten av materie som kan isoleres ved kjemiske metoder.

\vskip 10pt

Et {\bf grunnstoff} er et stoff sammensatt av atomer som alle har samme antall protoner i kjernen (f.eks. hydrogen, helium, nitrogen, jern, cesium, fluor).

\vskip 10pt

Et {\bf molekyl} er den minste enheten av materie som består av to eller flere atomer bundet sammen ved elektroninteraksjon i et fast forhold (f.eks. vann: H$_{2}$O) -- den minste enheten i en {\it forbindelse}.

\vskip 10pt

En {\bf forbindelse} er et stoff sammensatt av identiske molekyler (f.eks. rent vann).

\vskip 10pt

En {\bf blanding} er et stoff sammensatt av ulike atomer eller molekyler som {\it ikke} er elektronisk bundet til hverandre.

\vskip 10pt

\begin{itemize}
\item{} Ett gram helt rent vann: {\bf Forbindelse}
\item{} Et proton: {\bf Partikkel}
\item{} Oksygen: {\bf Atom}, {\bf Molekyl} eller {\bf Grunnstoff}, avhengig av kontekst.
\item{} Karbondioksid (CO$_{2}$): {\bf Molekyl} eller {\bf Forbindelse}, avhengig av kontekst.
\item{} En kvark: {\bf Partikkel}
\item{} Bensin: {\bf Blanding}
\item{} Luft: {\bf Blanding}
\item{} Plutonium: {\bf Atom} eller {\bf Grunnstoff}, avhengig av kontekst.
\end{itemize} 

\vskip 10pt

``Oksygen'' kan vise til ett enkelt atom, og klassifiseres da som atom. Oksygen forekommer vanligvis som molekyler (to oksygenatomer bundet sammen: O$_{2}$), og da klassifiseres en slik atompar-enhet som molekyl. Hvis ``oksygen'' brukes om stofftype, klassifiseres det som grunnstoff.

\vskip 10pt

Tilsvarende gjelder for karbondioksid og plutonium. Ett enkelt plutoniumatom er selvfølgelig et atom. En gruppe bestående av ett karbonatom (C) og to oksygenatomer (O$_{2}$) utgjør ett CO$_{2}$-molekyl. Hvis referansen gjelder stofftype, klassifiseres plutonium som ``grunnstoff'' og karbondioksid som ``forbindelse''.

\vskip 10pt

Bensin og luft er begge blandinger, der hvert stoff består av flere typer molekyler.

%(END_ANSWER)





%(BEGIN_NOTES)


%INDEX% Kjemi, grunnprinsipper: atom, partikkel, grunnstoff, molekyl, forbindelse, ion, blanding

%(END_NOTES)

